\section{Optimasi \textit{Simultaneous Perturbation Stochastic Approximation} (SPSA)}
\label{sec:4_6_optimasi_spsa}

\subsection{Dinamika Pembaruan Parameter SPSA}

Sebagai sebuah prosedur mitigasi atas rentannya observasi hambatan kinerja sistem pada sirkuit kuantum deterministik, pembaruan iteratif parameter matriks simulasi pada penelitian ini secara utuh diarahkan menuju penggunaan algoritma optimasi \textit{Simultaneous Perturbation Stochastic Approximation} (SPSA) \cite{spall_implementation_1998}. Algoritma turunan stokastik ini diinisialisasi secara khusus sebagai perangkat fungsional untuk menggantikan metode \textit{Gradient Descent} konvensional yang kerap kali mengalami kegagalan saat mengeksekusi navigasi pada ruang fase portofolio dimensi tinggi. Mekanika komputasi SPSA sangat diunggulkan secara operasional karena algoritma ini terbukti secara empiris sanggup menyelesaikan evaluasi fungsi gradien untuk keseluruhan komponen arsitektur sirkuit matriks energi secara komprehensif. Reduksi drastis pada kuantitas evaluasi matriks fungsional tersebut membuat implementasi algoritma simulasi berbasis SPSA ini pada akhirnya sukses menetralkan lonjakan beban komputasional yang mengganggu presisi fungsional iterasi algoritma \cite{nielsen_quantum_2010}.

Meskipun secara fundamental algoritma matriks ini memanipulasi turunan fungsi objektif melalui pendekatan aproksimasi stokastik yang tidak sempurna, parameter hasil observasi sirkuit VQE akan tetap dieksekusi secara stabil menuju konvergensi nilai rotasi optimal secara global. Keunikan probabilitas peluruhan parameter ini terjadi secara langsung akibat polaritas alamiah dari distribusi turunan fungsi Rademacher yang secara fungsional mampu mendegradasi sisa galat aproksimasi pada dimensi perhitungan operasi berpangkat genap. Oleh sebab itu, tingkat ekspektasi komputasional turunan energi pada sirkuit kuantum secara rasional diyakini bernilai sangat identik dengan besaran matriks evaluasi gradien riil yang menyusun lanskap portofolio fungsional yang sesungguhnya \cite{spall_implementation_1998}. Atas dasar pembuktian ekuivalensi komprehensif antara tingkat ketelitian algoritma SPSA dan efisiensi waktu operasinya, prosedur uji coba observasi parameter portofolio pada sirkuit riil \textit{hardware} NISQ dapat dieksekusi secara aplikatif \cite{nielsen_quantum_2010}.

\subsection{Implementasi Numerik dan Konvergensi Energi Optimal}
\label{subsec:4_6_2_numerik_fisis}

Guna mentranslasikan kompleksitas matematis SPSA menjadi sebuah pemahaman operasional yang jauh lebih konkret bagi pembaca akademik, penjabaran satu siklus tahapan iterasi simulasi ini disajikan secara komprehensif. Proses manualisasi fungsi komputasi ini memanfaatkan formulasi bentuk murni Hamiltonian Ising berdimensi dua aset ($N=2$) yang secara utuh telah diturunkan pada pembahasan \S \ref{sec:4_3_pemetaan_hamiltonian} sebelumnya. Bentuk fungsional Hamiltonian Ising yang telah memperhitungkan faktor penalti pelanggaran anggaran tersebut diuraikan sebagai berikut:
\begin{equation}
    \hat{\mathcal{H}} \approx 0,1736 \hat{Z}_1 + 0,0932 \hat{Z}_2 + 2,5014 \hat{Z}_1 \hat{Z}_2 + 2,2320.
\end{equation}
Langkah-langkah struktural dari proses pembaruan parameter fungsional menggunakan metode SPSA ini selanjutnya dieksekusi secara mendetail melalui sekuens perhitungan sistematis yang meliputi inisialisasi gangguan hingga pada tahapan kalkulasi akhir estimator gradien.

\begin{enumerate}
    \item \textbf{Inisialisasi Parameter dan Gangguan:}
    Misalkan posisi parameter awal adalah $\theta^{(0)} = [\pi/2, \pi/2]^T \approx [1,5708, 1,5708]^T$. Dipilih koefisien gangguan $c_k = 0,1$ dan vektor gangguan Rademacher $\Delta_k = [1, -1]^T$. Maka, titik evaluasi parameter adalah:
    \begin{equation*}
        \theta_+ = \theta^{(0)} + c_k \Delta_k = [1,6708, 1,4708]^T, \quad \theta_- = \theta^{(0)} - c_k \Delta_k = [1,4708, 1,6708]^T
    \end{equation*}

    \item \textbf{Estimasi Nilai Ekspektasi Energi ($E_{\pm}$):}
    Guna menyederhanakan ilustrasi, kita asumsikan \textit{ansatz} sederhana di mana $\langle \hat{Z}_i \rangle = \cos(\theta_i)$. Maka fungsi biaya energinya menjadi:
    \begin{equation*}
        E(\theta) = 0,1736 \cos(\theta_1) + 0,0932 \cos(\theta_2) + 2,5014 \cos(\theta_1)\cos(\theta_2) + 2,2320
    \end{equation*}
    Substitusi nilai $\theta_+$ dan $\theta_-$ menghasilkan:
    \begin{itemize}
        \item $E_+ \approx 0,1736 \cos(1,6708) + 0,0932 \cos(1,4708) + 2,5014 \cos(1,6708)\cos(1,4708) + 2,2320 \approx 2,1990$.
        \item $E_- \approx 0,1736 \cos(1,4708) + 0,0932 \cos(1,6708) + 2,5014 \cos(1,4708)\cos(1,6708) + 2,2320 \approx 2,2151$.
    \end{itemize}

    \item \textbf{Kalkulasi Estimator Gradien ($\hat{g}$):}
    Selisih energi $\Delta E = E_+ - E_- = 2,1990 - 2,2151 = -0,0161$. Estimator gradien dihitung sebagai:
    \begin{equation*}
        \hat{g} = \frac{\Delta E}{2 c_k \Delta_k} = \begin{pmatrix} -0,0161 / (2 \cdot 0,1 \cdot 1) \\ -0,0161 / (2 \cdot 0,1 \cdot -1) \end{pmatrix} = \begin{pmatrix} -0,0805 \\ 0,0805 \end{pmatrix}
    \end{equation*}

    \item \textbf{Pembaruan Parameter ($\theta^{(1)}$):}
    Dengan menggunakan \textit{learning rate} $\eta = 0,1$, parameter diperbarui menjadi:
    \begin{equation*}
        \theta^{(1)} = \theta^{(0)} - \eta \hat{g} = \begin{pmatrix} 1,5708 \\ 1,5708 \end{pmatrix} - 0,1 \begin{pmatrix} -0,0805 \\ 0,0805 \end{pmatrix} = \begin{pmatrix} 1,5789 \\ 1,5627 \end{pmatrix}
    \end{equation*}
\end{enumerate}

Melalui skema iterasi stohastik yang dijabarkan tersebut, serangkaian hasil simulasi pengujian proses optimasi SPSA pada jendela data kuantum di bulan Januari 2021 berhasil mendeteksi keadaan optimal bagi distribusi energi fungsional untuk setiap kemungkinan konfigurasi \textit{bitstring} secara presisi. Dari evaluasi nilai komputasi ini didapatkan hasil bahwa keadaan $|00\rangle$ berasosiasi dengan $E \approx 5,0$ yang mengindikasikan status terlarang, sedangkan konfigurasi $|01\rangle$ memiliki matriks energi $E \approx -0,189$ sebagai solusi sub-optimal. Sementara itu, solusi \textit{ground state} global fungsional yang merepresentasikan optimasi portofolio terbaik dipetakan secara absolut pada konfigurasi fisis status $|10\rangle$ dengan perolehan rekor energi ekspektasi bernilai $E \approx -0,350$. Untuk menunjang implementasi fungsi komputasi berulang ini secara otonom dalam skala analisis portofolio kuantum riil, kerangka dasar algoritma stokastik tersebut diakomodasi murni secara leksikal pada \textit{script} pemrograman di dalam Daftar Kode \ref{lst:run_spsa_test}.

\begin{lstlisting}[language=Python, caption={Implementasi algoritma dasar \textit{Simultaneous Perturbation Stochastic Approximation} (SPSA)}, label={lst:run_spsa_test}]
import numpy as np

def run_spsa_test(cost_circuit, n_params, init_params=None, a_base=0.1, c_base=0.1, max_iters=30, batch_size=25):
    """Versi ringan SPSA untuk pengetesan Learning Rate."""
    rng = np.random.default_rng(42)
    if init_params is not None:
        params = init_params.copy()
    else:
        params = rng.uniform(0, 2 * np.pi, n_params)
    
    A_coeff = max_iters * 0.1
    alpha_exp = 0.602
    gamma_exp = 0.101
    
    total_iters = 0
    while total_iters < max_iters:
        for _ in range(batch_size):
            k = total_iters
            a_k = a_base / (A_coeff + k + 1) ** alpha_exp
            c_k = c_base / (k + 1) ** gamma_exp
            
            # Formulasi parameter perturbasi serentak (Rademacher)
            delta = 2 * rng.integers(0, 2, size=n_params) - 1
            
            # Estimasi fungsi biaya energi (cost circuit)
            cost_plus = float(cost_circuit(params + c_k * delta))
            cost_minus = float(cost_circuit(params - c_k * delta))
            
            # Kalkulasi gradien
            grad = (cost_plus - cost_minus) / (2 * c_k * delta)
            params = params - a_k * grad
            
            total_iters += 1
            
    return float(cost_circuit(params))
\end{lstlisting}

